\begin{lemma}
\label{editorial-source-lemma-finite-finitely-presented-extension}
Let $R \to S$ be a finite and finitely presented ring map.
Let $M$ be an $S$-module.
Then $M$ is finitely presented as an $R$-module if and only if
$M$ is finitely presented as an $S$-module.
\end{lemma}

\begin{proof}
If $M$ is finitely presented over $R$, Lemma
\ref{algebra-lemma-finitely-presented-over-subring} applies: the given finite
map is of finite type, because its module generators also generate
its algebra. Hence $M$ is finitely presented over $S$.
To see the other assume that $M$ is finitely presented as an $S$-module.
Pick a presentation
$$
S^{\oplus m} \longrightarrow
S^{\oplus n} \longrightarrow
M \longrightarrow 0
$$
The kernel of $S^{\oplus n}\to M$ is the image of the displayed
$S^{\oplus m}$. A finite $R$-generating family of $S$ in each of
these $m$ summands maps to a finite $R$-generating family of that
kernel. Thus from
Lemma \ref{algebra-lemma-extension}
we see that it suffices to prove that $S$ is finitely presented as an
$R$-module.

\medskip\noindent
Pick $y_1, \ldots, y_n \in S$ such that $y_1, \ldots, y_n$ generate $S$
as an $R$-module. By Lemma \ref{editorial-source-lemma-characterize-integral-element}
each $y_i$ is integral over $R$. Choose monic polynomials
$P_i(x) \in R[x]$ with $P_i(y_i) = 0$. Consider the ring
$$
S' = R[x_1, \ldots, x_n]/(P_1(x_1), \ldots, P_n(x_n))
$$
Then we see that $S$ is of finite presentation as an $S'$-algebra
by Lemma \ref{algebra-lemma-compose-finite-type}. Since $S' \to S$ is surjective,
the kernel $J = \Ker(S' \to S)$ is finitely generated as an ideal by
Lemma \ref{algebra-lemma-finite-presentation-independent}. Hence $J$ is a finite
$S'$-module (immediate from the definitions).
Thus $S = \Coker(J \to S')$ is  of finite presentation as an $S'$-module
by Lemma \ref{algebra-lemma-extension}.
Hence, arguing as in the first paragraph, it suffices to show that $S'$ is
of finite presentation as an $R$-module. Actually, $S'$ is free as an
$R$-module with basis the monomials $x_1^{e_1} \ldots x_n^{e_n}$
for $0 \leq e_i < \deg(P_i)$. Namely, write $R \to S'$ as the composition
$$
R \to R[x_1]/(P_1(x_1)) \to R[x_1, x_2]/(P_1(x_1), P_2(x_2)) \to
\ldots \to S'
$$
This shows that the $i$th ring in this sequence is free as a module over the
$(i - 1)$st one with basis $1, x_i, \ldots, x_i^{\deg(P_i) - 1}$. The result
follows by induction, with the following division argument proving
both spanning and independence over every intermediate coefficient
ring. If $P_i(x_i)=x_i^{d_i}+\sum_{j<d_i}a_{ij}x_i^j$ is monic,
subtract the leading coefficient times $x_i^{e-d_i}P_i$ from
a polynomial of degree $e\geq d_i$. Repeating decreases the degree
and leaves a remainder of degree less than $d_i$, retaining each
coefficient $a_{ij}$. Uniqueness holds because for nonzero $Q$,
the product $P_iQ$ has leading coefficient equal to that of $Q$
and degree $d_i+\deg Q$, even if the coefficient ring has
zero divisors. Thus no nonzero remainder of degree less than
$d_i$ belongs to $(P_i)$. Choosing $d_i\geq1$ is always possible;
a monic equation of degree zero can be multiplied by $x_i$.
Applying the unique remainder successively in $x_1,\ldots,x_n$
gives exactly the displayed monomial basis of $S'$ over $R$,
of size $D=\prod_i d_i$ (the empty product is $1$).

To spell out the final presentations, let $h_1,\ldots,h_a$
generate the original ideal $J\subset S'$. The $aD$ elements
$h_j x_1^{e_1}\cdots x_n^{e_n}$, with the displayed exponent
bounds, generate $J$ as an $R$-module. Indeed expand each
coefficient in an ideal expression $\sum_j c_jh_j$ uniquely
in that monomial basis. Their coordinate vectors give a map
$R^{\oplus aD}\to R^{\oplus D}$ with cokernel the original
$S=S'/J$. This is an explicit finite $R$-presentation of $S$.
For the original $S^{\oplus n}$ of the first paragraph, take
one copy of this presentation per summand. That number of summands
is independent of the number of chosen $y_i$ in this paragraph.
Lift the finite $R$-generating family
of the kernel there to this finite free $R$-module. The lifted
vectors together with its existing finitely many relation vectors
generate the full kernel onto $M$: subtracting the lifts from
any vector mapping to zero leaves an old relation. This supplies
a finite $R$-presentation of the same $M$ and finishes the proof.
\end{proof}